\chapter{Etikken og stemma}

\begin{motto}
Påstanden er ikkje at design manglar ei etikk.\\ Påstanden, rett stilt, er at etikken er rik — men uhandhevd.
\end{motto}

Her kjem forsvaret til sitt vanskelegaste kapittel, og eg vil ikkje vinne det med eit triks. \textsc{Grunnlaust} hevdar at design manglar ei eiga etikk, eit internasjonalt etisk organ, og ei stemme; at det ikkje finst nokon standard å bli utelukka frå. Den fyrste delen av denne påstanden er rett og slett feil, og det skal eg vise. Den andre delen inneheld eit reelt poeng som forsvaret må innrømme, ikkje bortforklare. Skiljet mellom dei to er heile saka, og eg trur den ærlege handsaminga av det styrkjer forsvaret meir enn ei triumferande ein ville gjort.

\section{Etikken finst, og ho er rik}

Lat oss byrje med det som er lett å verifisere, fordi det feller den sterke forma av påstanden med ein gong. Designfaget har ikkje éi etikk — det har mange, og fleire av dei er gjennomarbeidde, institusjonaliserte og levande. Det internasjonale designrådet (\textsc{ico-d}, tidlegare Icograda, grunnlagt 1963) har ein modellkodeks for profesjonell framferd som plasserer designaren med ansvar «for menneskeheita og ei berekraftig framtid», meint som mal for nasjonale foreiningar og for læreplanar \autocite{icod_code}. Verdsorganisasjonen for design (\textsc{wdo}, tidlegare \textsc{icsid}, grunnlagt 1957) har ein eigen etisk kodeks og eit komité for profesjonell praksis \autocite{wdo_ethics}. \textsc{Aiga}, verdas største designorganisasjon, har detaljerte standardar for profesjonell praksis som eksplisitt seier at ein designar bør avvise oppdrag som krev brot på dei etiske standardane \autocite{aiga_standards}; \textsc{idsa} bind medlemene sine til ein etisk kodeks som nemner miljø og framtidige generasjonar \autocite{idsa_ethics}.

Og dette er berre profesjonsorgana. Ved sida av dei ligg ein heil tradisjon av eksplisitt designetikk: First Things First-manifestet frå 1964, fornya i 2000, som kravde at designarar vende evnene sine mot det som tener samfunnet \autocite{firstthings2000}; Design Justice Network sine ti prinsipp frå 2018, med Sasha Costanza-Chock si utgreiing som akademisk grunnlag \autocite{costanzachock2020}; den skandinaviske tradisjonen for participatory design, der Pelle Ehn gjorde arbeidsplassdemokrati og utjamning av makt til eksplisitte designmål \autocite{ehn1992scandinavian}; og nyare arbeid som Cennydd Bowles si \textit{Future Ethics}, som omset moderne etisk teori til praktisk designarbeid \autocite{bowles2018future}. Den som seier at design manglar ei etikk, har ikkje sett etter. Faget har skrive om etikk i seksti år, frå fleire kantar, med stigande presisjon.

\section{Stemma finst òg}

Det same gjeld påstanden om at faget manglar ei stemme. I 2017 signerte tre internasjonale designorganisasjonar Montréal-deklarasjonen i nærvær av tre \textsc{fn}-organ, og slo fast designets ansvar for ei berekraftig, rettferdig og kulturelt mangfaldig verd \autocite{montreal2017}. Cumulus-foreininga, med over 350 kunst- og designhøgskular i 60 land, forplikta medlemene sine til eit globalt ansvar gjennom Kyoto-deklarasjonen \autocite{kyoto2008}. I Noreg er \textsc{doga} eit statleg kompetansesenter som gjer sosial, miljømessig og økonomisk verdi til vurderingskriterium, og Grafill ein landsdekkande fagorganisasjon som tek stilling til nye etiske spørsmål, frå opphavsrett til kunstig intelligens. Når designfaget vil tale, har det altså adresser å tale frå. At desse stemmene ikkje alltid blir høyrde så høgt som klimavitskapen sin, er eit spørsmål om makt og merksemd, ikkje om fråvær.

\section{Den ærlege innrømminga: handheving manglar}

No til det forsvaret ikkje skal pynte på, fordi det er sant. Sjå kva alle desse kodeksane og deklarasjonane har felles: nesten ingen av dei kan \emph{handhevast} over individet. \textsc{Aiga} og \textsc{idsa} kan i prinsippet ekskludere eit medlem, men eksklusjonen råkar berre medlemskapet, ikkje retten til å kalle seg designar eller ta oppdrag. Det finst inga lisensiering. Det finst ingen verna tittel. Og det finst ikkje noko internasjonalt etisk tilsynsorgan for design med makt til å frata ein utøvar retten til å praktisere.

Kontrasten med nabofaget er tydeleg og ubehageleg. Arkitektane har det design manglar: \textsc{riba} sin kodeks kan føre til privat åtvaring, offentleg refs, suspensjon eller eksklusjon gjennom eit formelt disiplinærpanel \autocite{riba2019code}; \textsc{aia} sitt etiske råd handhevar obligatoriske åtferdsreglar og publiserer straffene \autocite{aia_ethics}. Dette er kva det vil seie å ha ein standard ein kan bli utelukka frå. Design har det stort sett ikkje. På dette punktet har \textsc{Grunnlaust} rett, og forsvaret tener ingenting på å nekte for det.

\section{Kva innrømminga faktisk inneber}

Men sjå nøye på kva som følgjer, og kva som ikkje følgjer, av denne innrømminga. Det som følgjer, er at designetikken er \emph{uhandhevd} og \emph{ikkje-lisensiert}. Det som \emph{ikkje} følgjer, er at han ikkje finst, eller at faget manglar eit etisk fundament. Påstanden må reformulerast: ikkje «design har inga etikk», men «design har ein rik etikk utan tennene til å tvinge han igjennom». Det er ein heilt annan, og langt svakare, kritikk — og det er ein kritikk om institusjonell modning, ikkje om intellektuelt fundament.

Og modning tek tid på ein måte som er verdt å minne om. Legestanden fekk ikkje sin handhevbare etikk over natta; han voks fram over hundreår, gjennom profesjonskampar, lovgjeving og institusjonsbygging, lenge etter at den medisinske kunnskapen var etablert. At design — eit fag som fyrst vart akademisk for omtrent hundre år sidan, og som har vakse eksplosivt dei siste tiåra — enno ikkje har bygd det same handhevingsapparatet, er kva ein ventar av eit ungt fag, ikkje prov på at det er fundamentlaust. Etikken er der. Apparatet for å tvinge han igjennom er under bygging. Å forveksle det andre fråvêret med det fyrste er å forveksle ungdom med tomleik.

\vignett

Forsvaret har no ført fram kunnskapsforma, metoden, kumulasjonen, det vrange som innsikt, meininga, mennesket og etikken. Det står att å vise at desse ikkje berre er internasjonale abstraksjonar, men at dei er bygde, konkret, i eit fagmiljø — og her vender vi heim, til den norske og nordiske designforskinga, som har gjort meir enn dei fleste for å gje design eit fundament det kan stå på.

\vidarelesing{\fullcite{costanzachock2020}; \fullcite{ehn1992scandinavian}; \fullcite{icod_code}; \fullcite{aiga_standards}; \fullcite{riba2019code}; \fullcite{bowles2018future}.}
